\documentclass[a4paper,12pt]{article}

\usepackage[T1]{fontenc}
\usepackage[italian]{babel}
\usepackage{graphicx}
\usepackage{geometry}
\usepackage{tabularx}
\usepackage[table]{xcolor}
\usepackage[most]{tcolorbox}
\usepackage{fancyhdr}
\usepackage[hidelinks]{hyperref}

\newcommand{\groupmail}{alittlebyte19@gmail.com}
\newcommand{\grouplogo}{.github/site-src/assets/logo_alittlebyte.png}
\newcommand{\groupsymbol}{.github/site-src/assets/solo_logo.png}

\geometry{
    left=2.5cm,
    right=2.5cm,
    top=2.5cm,
    bottom=2.5cm,
    headheight=331pt,
    includeheadfoot
}

\pagestyle{fancy}
\fancyhf{}

\renewcommand{\headrulewidth}{0pt}
\fancyhead{}

\renewcommand{\footrulewidth}{0.4pt}
\fancyfoot[L]{\includegraphics[height=0.6cm]{\groupsymbol}\hspace{0.45em}aLittleByte}
\fancyfoot[R]{Norme di Progetto}

\fancypagestyle{firstpage}{
    \fancyhf{}
    \renewcommand{\headrulewidth}{0pt}
    \renewcommand{\footrulewidth}{0.4pt}
    \fancyhead[C]{%
        \makebox[\textwidth][c]{%
            \begin{tabular}{c}
                \includegraphics[height=10cm]{\grouplogo}\\[1ex]
                \small\texttt{\groupmail}
            \end{tabular}%
        }
    }
    \fancyfoot[L]{\includegraphics[height=0.6cm]{\groupsymbol}\hspace{0.45em}aLittleByte}
    \fancyfoot[R]{Norme di Progetto}
}

\begin{document}

\thispagestyle{firstpage}

\vspace*{0.5cm}

\begin{center}
    {\LARGE \textbf{Norme di Progetto}}
\end{center}

\vspace{1.5cm}

\begin{center}
    \renewcommand{\arraystretch}{1.3}
    \begin{tcolorbox}[
        width=0.92\textwidth,
        colback=gray!6,
        colframe=gray!45,
        boxrule=0.5pt,
        arc=2mm,
        boxsep=0pt,
        left=8pt, right=8pt, top=8pt, bottom=8pt
    ]
        \begin{tabularx}{\linewidth}{>{\bfseries}p{3.5cm} X}
            Gruppo      & aLittleByte (gruppo 19) \\
            Redattori   & Eleonora Bellet, Diego Belluz \\
            Verificatori& Leonardo Soligo, Angelica Gastaldello\\
            Destinatari & Prof. Tullio Vardanega, Prof. Riccardo Cardin \\
            Data        & -/-/2026 \\
        \end{tabularx}
    \end{tcolorbox}
\end{center}

\newpage

\newgeometry{left=2.5cm,right=2.5cm,top=2.5cm,bottom=2.5cm,includefoot}

\begin{center}
    {\Large \textbf{Registro delle Modifiche}}
\end{center}

\vspace{0.35cm}

\renewcommand{\arraystretch}{1.35}
\rowcolors{2}{gray!8}{white}
\begin{center}
    {\footnotesize
    \begin{tabularx}{\textwidth}{|p{1.4cm}|p{1.8cm}|X|p{2.4cm}|p{2.3cm}|p{2.3cm}|}
        \hline
        \rowcolor{gray!20}
        \textbf{Versione} & \textbf{Data} & \textbf{Descrizione} & \textbf{Redattore} & \textbf{Verificatore} & \textbf{Approvatore} \\
        \hline
        0.3.0 & 4/04/2026 & Architettura dei processi & Diego Belluz & Eleonora Bellet & - \\
        \hline
        0.2.0 & 31/03/2026 & Scopo del documento e scopo del prodotto & Eleonora Bellet & Diego Belluz & - \\
        \hline
        0.1.0 & 31/03/2026 & Struttura del documento & Eleonora Bellet & Diego Belluz & - \\
        \hline
        
    \end{tabularx}
    }
    
\end{center}
\newpage

\begin{center} 
    {\Large \textbf{Indice}}
\end{center}

\vspace{0.6cm}

\tableofcontents

\newpage
\pagenumbering{arabic}
\setcounter{page}{4}
\fancyfoot[R]{Norme di Progetto\hspace{0.8em}\thepage}

\section{Introduzione}

\subsection{Scopo del documento}
Il presente documento ha l'obiettivo di definire e regolamentare il \textit{Way of Working (WoW)} adottato dal gruppo \textit{ALittleByte} per la realizzazione del progetto didattico. Al suo interno sono specificate le procedure operative, gli strumenti adottati e le regole di collaborazione interne al team, al fine di garantire un approccio strutturato, efficiente e tracciabile per l'intero ciclo di vita del progetto Nexum. 

La stesura del seguente documento segue un approccio incrementale: le norme qui descritte verranno costantemente aggiornate e affinate nel tempo, supportando il miglioramento continuo dei processi organizzativi, di sviluppo e di supporto del gruppo.


\subsection{Scopo del prodotto}
Il prodotto da realizzare consiste in un'estensione innovativa per \textit{''Nexum"}, una piattaforma digitale sviluppata dall'azienda proponente \textit{Eggon}, dedicata alla comunicazione interna e alla gestione HR.
Nello specifico, il progetto mira a potenziare l'ecosistema Nexum attraverso lo sviluppo e l'integrazione di due moduli basati sull'Intelligenza Artificiale:
\begin{itemize}
    \item \textbf{AI Assistant Generativo:} Un modulo dedicato agli utenti amministrativi (dashboard) che permette la creazione autonoma di contenuti aziendali (titolo, testo, immagini di copertina) a partire da un prompt, garantendo l'adeguamento automatico del tono e dello stile del messaggio.
    \item \textbf{AI Co-Pilot per Consulenti del Lavoro (CdL):} Un sistema avanzato in grado di riconoscere automaticamente la tipologia di documenti caricati (es. cedolini, comunicazioni), estrarne i destinatari tramite tecniche di OCR ed \textbf{Entity Resolution}, e automatizzare lo smistamento e il dispaccio massivo.
\end{itemize}

L'obiettivo finale è migliorare l'efficienza dei processi HR e semplificare l'interoperabilità tra le aziende e gli studi dei Consulenti del Lavoro. Per fare ciò il gruppo ALittleByte si impegna a completare il progetto entro la data fissata 10 Settembre 2026, con un costo totale di budget di 12.575 €. 

Per ulteriori informazioni si rimanda al documento Analisi dei Requisiti. 


\subsection{Glossario}

Al fine di prevenire qualsiasi ambiguità linguistica o concettuale tra i membri del team e gli stakeholder, è stato redatto un documento separato denominato Glossario. Esso contiene la spiegazione puntuale di tutti i termini tecnici, di dominio (es. LLM, OCR, Dispaccio, ecc.) e degli acronimi utilizzati all'interno della documentazione di progetto.


\subsection{Riferimenti}

\subsubsection{Normativi}

\subsubsection{Informativi}

\section{Architettura dei processi}
Per la realizzazione di un prodotto di qualità che soddisfi le esigenze del cliente e degli utenti, il gruppo aLittleByte ha deciso di adottare un'architettura specifica e suddividere i vari processi di costruzione del software in tre macro-aree:
\begin{itemize}
    \item Processi primari
    \item Processi di supporto
    \item Processi organizzativi
\end{itemize}

\subsection{Processi Primari}
In questa sezione viene descritto come nasce il software. I processi primari indicano tutte quelle attività che sono direttamente coinvolte nella creazione del prodotto software, è possibile classificare i processi primari in:
\begin{itemize}
    \item Processi di fornitura 
    \item Processi di sviluppo
\end{itemize}

\subsection{Processi di fornitura}
I processi di fornitura comprendono tutte le attività preliminari che permettono di avviare lo sviluppo del progetto, in particolar modo la pianificazione: suddivisione del lavoro in fasi e milestone come rtb e pb, l’analisi: come vengono analizzati i capitolati, la stima: come si stimano vari costi e stime.


\end{document}
